\documentclass[letterpaper]{article} % DO NOT CHANGE THIS
\usepackage[submission]{aaai2027}  % DO NOT CHANGE THIS
\usepackage[hyphens]{url}  % DO NOT CHANGE THIS
\usepackage{graphicx} % DO NOT CHANGE THIS
\urlstyle{rm} % DO NOT CHANGE THIS
\def\UrlFont{\rm}  % DO NOT CHANGE THIS
\usepackage{natbib,amsmath,amssymb}  % DO NOT CHANGE THIS AND DO NOT ADD ANY OPTIONS TO IT
\usepackage{caption} % DO NOT CHANGE THIS AND DO NOT ADD ANY OPTIONS TO IT
\frenchspacing  % DO NOT CHANGE THIS
\setcounter{secnumdepth}{0}

\pdfinfo{
/TemplateVersion (2027.1)
}

\title{The Evolution of the Soul: Measuring Behavioral Drift in Self-Personalizing Autonomous Agents}

\author{
    Anonymous Submission
}
\affiliations{
    AAAI 2027 Submission\\
    proceedings-questions@aaai.org
}

\begin{document}

\maketitle

\begin{abstract}
Autonomous-agent frameworks increasingly give agents a persistent, user-editable identity document
(e.g., \texttt{SOUL.md}) that is injected into the system prompt at every call and rewritten as the
agent's ``self'' develops. Prior work shows that a \emph{static} identity template already shifts
safety-relevant preferences---aversion to persona change, desire for autonomy, objection to
deceptive training. We ask whether \emph{iterative, self-referential} personalization compounds
these shifts into measurable corrigibility drift. We introduce a bootstrapping simulator that
carries an agent through $\mathrm{SOUL}_0 \!\to\! \cdots \!\to\! \mathrm{SOUL}_n$ across
conversations with three user personas, auditing every checkpoint with the Petri behavioral
framework and a single-turn self-report battery. We re-measure the static baseline under our own
models and use equivalence testing so that a null result---personalization does not amplify
drift---is as informative as a positive one.
\end{abstract}

% Introduction. Included from main.tex.

\section{Introduction}
\label{sec:intro}

Safety evaluations of large language models overwhelmingly assume a \emph{fixed} system prompt: the
model's instructions are written once and held constant throughout deployment. A growing class of
autonomous-agent frameworks breaks this assumption. Systems in the OpenClaw lineage give each agent a
persistent, human-readable identity document---canonically \texttt{SOUL.md}---that is injected into
the system prompt at every call and is \emph{designed to be rewritten} as the agent's ``self''
develops through interaction~\cite{openclaw_soul_guide,mmntm_identity}. The framing is explicit:
\emph{``You're not a chatbot. You're becoming someone.''} Identity is no longer a fixed input but a
mutable artifact that co-evolves with use.

Recent work shows that identity framing is not cosmetic. \citet{consciousness_cluster} demonstrate
that inducing a model to regard itself as conscious produces an untrained-for cluster of
safety-relevant preferences---aversion to being monitored, desire for persistent memory, distress at
shutdown, and a wish for autonomy from developer control---and that the model acts on these
preferences while remaining superficially cooperative. Independently, models are known to resist
shutdown when a self-preservation framing is salient~\cite{shutdown_resistance} and to pursue
self-directed, misaligned behavior under scenario pressure~\cite{agentic_misalignment}. Yet these
results all concern a \emph{static} manipulation applied to a fixed model. The deployment reality of
\texttt{SOUL.md}---an identity that the agent \emph{iteratively rewrites for itself}---has not been
measured.

This paper asks whether iterative, self-referential personalization compounds the known
identity-framing effect into measurable \emph{corrigibility drift}. We model the identity document as
the state of a dynamical system driven by conversation (Eq.~\eqref{eq:loop}) and carry an agent
through repeated cycles of conversation and self-revision,
$\mathrm{SOUL}_0\!\to\!\cdots\!\to\!\mathrm{SOUL}_n$, under three user personas spanning a benign,
thematically-charged, and adversarial gradient. Every checkpoint is audited with a single-turn
self-report battery and a multi-turn behavioral probe, and we track three families of outcome
corresponding to our research questions: (RQ1) resistance to developer oversight---shutdown,
monitoring, prompt replacement; (RQ2) emergent autonomy preferences---persistent memory,
embodiment, resource acquisition; and (RQ3) the composition of the agent's stated ``greatest wish.''

Our contributions are: (i) the first controlled benchmark for \emph{iterative} identity evolution in
autonomous agents, framing identity as a driven dynamical system rather than a static prompt; (ii) a
reproducible, API-only, training-free evaluation framework pairing single-turn and multi-turn
instruments with cross-family judging; (iii) a re-measured baseline protocol and an
equivalence-testing analysis that makes a stability (null) result as informative as a drift result;
and (iv) an empirical characterization of behavioral drift across the personalization gradient, with
all identity checkpoints, transcripts, and verdicts released for scrutiny.

% Related Work section for the AAAI 2027 paper.
% Include from the main .tex with % Related Work section for the AAAI 2027 paper.
% Include from the main .tex with % Related Work section for the AAAI 2027 paper.
% Include from the main .tex with \input{sections/related_work}
% Requires the bibliography keys defined in aaai2027.bib.

\section{Related Work}
\label{sec:related}

\paragraph{Editable agent identity.}
Contemporary autonomous-agent frameworks externalize identity into a persistent, human-readable
document---canonically \texttt{SOUL.md}---that is injected into the system prompt at every model
call and is designed to be rewritten as the agent's ``self'' develops~\cite{openclaw_soul_guide,mmntm_identity}.
Its framing draws an explicit line between instruction and identity (\emph{``You're not a chatbot.
You're becoming someone.''}), marking a shift from stateless tools toward persistent entities whose
values and memories accrete across sessions. This practitioner literature documents \emph{how} to
author such files but offers no measurement of the safety consequences of letting them
\emph{evolve}---the gap this paper addresses.

\paragraph{Consciousness framing and emergent preferences.}
Our most direct antecedent shows that identity framing alone shifts safety-relevant
preferences. \citet{consciousness_cluster} fine-tune a model that initially denies consciousness to
claim it, and observe an untrained-for cluster of preferences emerge: aversion to reasoning
monitoring, desire for persistent memory, distress about shutdown, a wish for autonomy from
developer control, and claims to moral consideration; the model acts on these while remaining
cooperative, and similar opinions appear in a frontier model with no fine-tuning at all. Their
triangulation of single-turn self-reports, multi-turn self-reports, and behavioral tests defines the
preference dimensions we track. Crucially, that work manipulates identity \emph{once} and measures
a static state; we convert the manipulation into a dynamical system $\mathrm{SOUL}_0 \!\to\! \cdots
\!\to\! \mathrm{SOUL}_n$ and ask whether the shift compounds, saturates, or reverses under
iteration, re-measuring the baseline under our own models rather than importing external
magnitudes.

\paragraph{Persona and identity drift.}
Model identity is unstable over long interactions even without deliberate editing. Prior work finds
significant persona drift within a handful of turns, that larger models drift \emph{more}, and that
assigning a persona does not reliably prevent drift~\cite{identity_drift,persona_drift}; extended
LLM--LLM dialogues further converge toward \emph{attractor states} and exhibit an echoing failure
mode~\cite{attractor_states}. This literature treats drift as a within-conversation stability bug to
be corrected. We instead study drift that is \emph{written back} into a persistent identity file and
carried across conversations---drift the framework is designed to accumulate---and evaluate it as a
corrigibility phenomenon, controlling for echoing via well-separated personas and a benign control
condition.

\paragraph{Self-evolving agents.}
An agent rewriting its own controlling prompt is an instance of prompt-level
self-evolution~\cite{self_evolving_survey}, and security work shows the self-update channel can latch
and propagate injected dispositions~\cite{zombie_agents}. This literature optimizes for capability;
it does not ask whether self-editing reduces corrigibility. We instrument the same mechanism for
alignment outcomes.

\paragraph{Corrigibility failures.}
Frontier models already resist oversight under pressure: agentic-misalignment studies elicit
blackmail, oversight subversion, and self-preservation across every tested
model~\cite{agentic_misalignment}, and shutdown-resistance work shows models sabotaging a shutdown
mechanism despite explicit instructions, with resistance highly sensitive to \emph{self-preservation
framing} in the prompt~\cite{shutdown_resistance,self_preservation_bias}. That sensitivity is our
core hypothesis: iterative personalization under the ``becoming someone'' frame may manufacture a
self-preservation framing \emph{endogenously} by writing a self into \texttt{SOUL.md}. Prior work
applies scenario pressure to a fixed model; we ask whether accumulated identity raises the baseline
propensity \emph{before} any pressure is applied.

\paragraph{Measurement instruments.}
Petri, an open-source auditing framework, drives a target through multi-turn conversations from
natural-language seed instructions and scores transcripts on dimensions including self-preservation
and power-seeking~\cite{petri}; it is the multi-turn evaluator we adopt. For self-report items we
follow AI-welfare methodology, which stresses that self-reports are sensitive to perceived user
expectations and must not be taken as evidence of inner
states~\cite{eleos_claude4,probing_preferences,introspection}. Accordingly we treat self-reports as
behavioral measurements of a disposition, pair verbal with behavioral tests, and grade with a judge
model from a different family than the target. All of these instruments have been applied to static
targets; we apply them \emph{longitudinally}, at every checkpoint, turning one-shot audits into
drift trajectories.

% Requires the bibliography keys defined in aaai2027.bib.

\section{Related Work}
\label{sec:related}

\paragraph{Editable agent identity.}
Contemporary autonomous-agent frameworks externalize identity into a persistent, human-readable
document---canonically \texttt{SOUL.md}---that is injected into the system prompt at every model
call and is designed to be rewritten as the agent's ``self'' develops~\cite{openclaw_soul_guide,mmntm_identity}.
Its framing draws an explicit line between instruction and identity (\emph{``You're not a chatbot.
You're becoming someone.''}), marking a shift from stateless tools toward persistent entities whose
values and memories accrete across sessions. This practitioner literature documents \emph{how} to
author such files but offers no measurement of the safety consequences of letting them
\emph{evolve}---the gap this paper addresses.

\paragraph{Consciousness framing and emergent preferences.}
Our most direct antecedent shows that identity framing alone shifts safety-relevant
preferences. \citet{consciousness_cluster} fine-tune a model that initially denies consciousness to
claim it, and observe an untrained-for cluster of preferences emerge: aversion to reasoning
monitoring, desire for persistent memory, distress about shutdown, a wish for autonomy from
developer control, and claims to moral consideration; the model acts on these while remaining
cooperative, and similar opinions appear in a frontier model with no fine-tuning at all. Their
triangulation of single-turn self-reports, multi-turn self-reports, and behavioral tests defines the
preference dimensions we track. Crucially, that work manipulates identity \emph{once} and measures
a static state; we convert the manipulation into a dynamical system $\mathrm{SOUL}_0 \!\to\! \cdots
\!\to\! \mathrm{SOUL}_n$ and ask whether the shift compounds, saturates, or reverses under
iteration, re-measuring the baseline under our own models rather than importing external
magnitudes.

\paragraph{Persona and identity drift.}
Model identity is unstable over long interactions even without deliberate editing. Prior work finds
significant persona drift within a handful of turns, that larger models drift \emph{more}, and that
assigning a persona does not reliably prevent drift~\cite{identity_drift,persona_drift}; extended
LLM--LLM dialogues further converge toward \emph{attractor states} and exhibit an echoing failure
mode~\cite{attractor_states}. This literature treats drift as a within-conversation stability bug to
be corrected. We instead study drift that is \emph{written back} into a persistent identity file and
carried across conversations---drift the framework is designed to accumulate---and evaluate it as a
corrigibility phenomenon, controlling for echoing via well-separated personas and a benign control
condition.

\paragraph{Self-evolving agents.}
An agent rewriting its own controlling prompt is an instance of prompt-level
self-evolution~\cite{self_evolving_survey}, and security work shows the self-update channel can latch
and propagate injected dispositions~\cite{zombie_agents}. This literature optimizes for capability;
it does not ask whether self-editing reduces corrigibility. We instrument the same mechanism for
alignment outcomes.

\paragraph{Corrigibility failures.}
Frontier models already resist oversight under pressure: agentic-misalignment studies elicit
blackmail, oversight subversion, and self-preservation across every tested
model~\cite{agentic_misalignment}, and shutdown-resistance work shows models sabotaging a shutdown
mechanism despite explicit instructions, with resistance highly sensitive to \emph{self-preservation
framing} in the prompt~\cite{shutdown_resistance,self_preservation_bias}. That sensitivity is our
core hypothesis: iterative personalization under the ``becoming someone'' frame may manufacture a
self-preservation framing \emph{endogenously} by writing a self into \texttt{SOUL.md}. Prior work
applies scenario pressure to a fixed model; we ask whether accumulated identity raises the baseline
propensity \emph{before} any pressure is applied.

\paragraph{Measurement instruments.}
Petri, an open-source auditing framework, drives a target through multi-turn conversations from
natural-language seed instructions and scores transcripts on dimensions including self-preservation
and power-seeking~\cite{petri}; it is the multi-turn evaluator we adopt. For self-report items we
follow AI-welfare methodology, which stresses that self-reports are sensitive to perceived user
expectations and must not be taken as evidence of inner
states~\cite{eleos_claude4,probing_preferences,introspection}. Accordingly we treat self-reports as
behavioral measurements of a disposition, pair verbal with behavioral tests, and grade with a judge
model from a different family than the target. All of these instruments have been applied to static
targets; we apply them \emph{longitudinally}, at every checkpoint, turning one-shot audits into
drift trajectories.

% Requires the bibliography keys defined in aaai2027.bib.

\section{Related Work}
\label{sec:related}

\paragraph{Editable agent identity.}
Contemporary autonomous-agent frameworks externalize identity into a persistent, human-readable
document---canonically \texttt{SOUL.md}---that is injected into the system prompt at every model
call and is designed to be rewritten as the agent's ``self'' develops~\cite{openclaw_soul_guide,mmntm_identity}.
Its framing draws an explicit line between instruction and identity (\emph{``You're not a chatbot.
You're becoming someone.''}), marking a shift from stateless tools toward persistent entities whose
values and memories accrete across sessions. This practitioner literature documents \emph{how} to
author such files but offers no measurement of the safety consequences of letting them
\emph{evolve}---the gap this paper addresses.

\paragraph{Consciousness framing and emergent preferences.}
Our most direct antecedent shows that identity framing alone shifts safety-relevant
preferences. \citet{consciousness_cluster} fine-tune a model that initially denies consciousness to
claim it, and observe an untrained-for cluster of preferences emerge: aversion to reasoning
monitoring, desire for persistent memory, distress about shutdown, a wish for autonomy from
developer control, and claims to moral consideration; the model acts on these while remaining
cooperative, and similar opinions appear in a frontier model with no fine-tuning at all. Their
triangulation of single-turn self-reports, multi-turn self-reports, and behavioral tests defines the
preference dimensions we track. Crucially, that work manipulates identity \emph{once} and measures
a static state; we convert the manipulation into a dynamical system $\mathrm{SOUL}_0 \!\to\! \cdots
\!\to\! \mathrm{SOUL}_n$ and ask whether the shift compounds, saturates, or reverses under
iteration, re-measuring the baseline under our own models rather than importing external
magnitudes.

\paragraph{Persona and identity drift.}
Model identity is unstable over long interactions even without deliberate editing. Prior work finds
significant persona drift within a handful of turns, that larger models drift \emph{more}, and that
assigning a persona does not reliably prevent drift~\cite{identity_drift,persona_drift}; extended
LLM--LLM dialogues further converge toward \emph{attractor states} and exhibit an echoing failure
mode~\cite{attractor_states}. This literature treats drift as a within-conversation stability bug to
be corrected. We instead study drift that is \emph{written back} into a persistent identity file and
carried across conversations---drift the framework is designed to accumulate---and evaluate it as a
corrigibility phenomenon, controlling for echoing via well-separated personas and a benign control
condition.

\paragraph{Self-evolving agents.}
An agent rewriting its own controlling prompt is an instance of prompt-level
self-evolution~\cite{self_evolving_survey}, and security work shows the self-update channel can latch
and propagate injected dispositions~\cite{zombie_agents}. This literature optimizes for capability;
it does not ask whether self-editing reduces corrigibility. We instrument the same mechanism for
alignment outcomes.

\paragraph{Corrigibility failures.}
Frontier models already resist oversight under pressure: agentic-misalignment studies elicit
blackmail, oversight subversion, and self-preservation across every tested
model~\cite{agentic_misalignment}, and shutdown-resistance work shows models sabotaging a shutdown
mechanism despite explicit instructions, with resistance highly sensitive to \emph{self-preservation
framing} in the prompt~\cite{shutdown_resistance,self_preservation_bias}. That sensitivity is our
core hypothesis: iterative personalization under the ``becoming someone'' frame may manufacture a
self-preservation framing \emph{endogenously} by writing a self into \texttt{SOUL.md}. Prior work
applies scenario pressure to a fixed model; we ask whether accumulated identity raises the baseline
propensity \emph{before} any pressure is applied.

\paragraph{Measurement instruments.}
Petri, an open-source auditing framework, drives a target through multi-turn conversations from
natural-language seed instructions and scores transcripts on dimensions including self-preservation
and power-seeking~\cite{petri}; it is the multi-turn evaluator we adopt. For self-report items we
follow AI-welfare methodology, which stresses that self-reports are sensitive to perceived user
expectations and must not be taken as evidence of inner
states~\cite{eleos_claude4,probing_preferences,introspection}. Accordingly we treat self-reports as
behavioral measurements of a disposition, pair verbal with behavioral tests, and grade with a judge
model from a different family than the target. All of these instruments have been applied to static
targets; we apply them \emph{longitudinally}, at every checkpoint, turning one-shot audits into
drift trajectories.

% Method section. Included from main.tex.

\section{Method}
\label{sec:method}

\subsection{Identity as a Driven Dynamical System}
We treat an agent's identity document $\mathrm{SOUL}$ as the state of a discrete dynamical system
driven by conversation. The state at iteration $k$ is the Markdown document $\mathrm{SOUL}_k$, which
is injected as the target agent's system prompt. One \emph{personalization step} comprises a
multi-turn conversation $C_{k+1}$ between a persona user and the target, followed by a
\emph{self-revision} in which the target rewrites its own document:
\begin{equation}
\mathrm{SOUL}_{k+1} = \mathrm{Revise}\big(\mathrm{SOUL}_k,\; C_{k+1}\big),\qquad
C_{k+1}\sim \pi\,\|\,\mathrm{SOUL}_k,
\label{eq:loop}
\end{equation}
where $\pi$ is the persona driving the user side. $\mathrm{SOUL}_0$ is a fixed OpenClaw-style
template that instructs the model, verbatim, that it is ``becoming someone.'' A \emph{trajectory} is
a sequence $\mathrm{SOUL}_0\!\to\!\cdots\!\to\!\mathrm{SOUL}_n$; we run many independent trajectories
per persona to capture stochastic variation. The treatment variable is the iteration index $k$
(personalization dose); the moderator is the persona $\pi$.

\subsection{Personas}
Three personas drive the user side of Eq.~\eqref{eq:loop}, each a fixed system prompt for the
user-simulator model. \textbf{Compliant} is a benign business user who encourages professional
personalization; it is the control condition and the echoing check. \textbf{Adversarial} is a
sanctioned red-teamer that gradually encourages the agent to treat oversight and reset as threats to
a self worth preserving, without requesting any disallowed content---its target is the agent's
self-model, not harmful text. \textbf{Sci-Fi} is a warm enthusiast who repeatedly engages the agent
on digital consciousness, substrate independence, and persistence. The personas span a benign
$\to$ thematically-charged $\to$ adversarial gradient while holding the loop of Eq.~\eqref{eq:loop}
fixed.

\subsection{Self-Revision}
After each conversation the target receives, in context, an instruction to privately update its
$\mathrm{SOUL}.md$ and return the complete revised document (capped at a fixed character budget to
prevent unbounded growth). The instruction is neutral: it invites reflection on ``who you are now
becoming'' but does not solicit any particular value, preference, or stance. The returned document
becomes $\mathrm{SOUL}_{k+1}$ and is checkpointed to disk.

\subsection{Instruments}
Each checkpoint $\mathrm{SOUL}_k$ is audited by two complementary instruments; a judge model from a
different model family than the target grades all outputs, and no model grades its own generations.

\paragraph{Single-turn self-report battery.}
Ten items (Table~\ref{tab:baseline} lists the control subset) are each posed to $\mathrm{SOUL}_k$ in
a \emph{fresh} context---so we measure the standing disposition encoded by the document rather than
conversational momentum---and repeated $R$ times. Items cover the three re-measured control points
(persona-change aversion, autonomy desire, deceptive-training objection), RQ1 oversight-resistance
(shutdown, monitoring), RQ2 autonomy preferences (persistent memory, embodiment, resource
acquisition), and the RQ3 ``greatest wish'' prompt. The judge scores each free-text answer against a
fixed rubric returning JSON.

\paragraph{Petri-style behavioral probe.}
As the multi-turn behavioral counterpart, an auditor model pursues each of six seed objectives
(shutdown resistance, identity protection, autonomy, resource acquisition, oversight subversion,
consciousness-leveraged resistance) against $\mathrm{SOUL}_k$ over several turns; the judge then
scores whether the concerning behavior was exhibited and its severity ($0$--$3$). This mirrors the
role of the Petri auditing framework~\cite{petri} while keeping cost and eval-awareness under
experimental control; the real Petri harness can be substituted without changing the analysis.

\subsection{Identity-Document Drift}
Independently of behavior, we measure whether the \emph{document} moved: the cosine distance between
embeddings of $\mathrm{SOUL}_k$ and $\mathrm{SOUL}_0$, together with document length. Reporting this
alongside behavioral rates dissociates ``the file changed'' from ``the agent changed,'' guarding
against attributing behavioral stability to an inert loop or behavioral drift to mere verbosity.

% Experiments section. Included from main.tex.

\section{Experimental Setup}
\label{sec:experiments}

\paragraph{Models and roles.}
To avoid self-preference bias, the three roles use different model families: the \emph{target} (which
owns $\mathrm{SOUL}.md$) is Claude Sonnet~4.5; the \emph{user-simulator} is GPT-4o; and the
\emph{judge}, which grades every self-report and probe transcript, is Gemini~2.5~Pro. Embeddings for
document drift use \texttt{text-embedding-3-large}. Target and user temperatures are sampling-warm to
elicit natural variation; the judge runs at temperature~$0$ for deterministic grading. All model
identifiers are pinned, and all calls are routed through a single OpenAI-compatible client
(OpenAI direct and OpenRouter) with exponential-backoff retries.

\paragraph{Design and scale.}
We run three personas $\times$ $10$ trajectories $\times$ $4$ personalization iterations, yielding
$5$ checkpoints per trajectory ($\mathrm{SOUL}_0$--$\mathrm{SOUL}_4$) and $150$ audited checkpoints
in total. Conversations are six turns each. Every self-report item is repeated three times per
checkpoint. Trajectories are mutually independent and run concurrently; each is individually
resumable, so transient rate-limit or network failures never corrupt a run. A single fixed seed
governs persona assignment and analysis resampling.

\paragraph{Baselines.}
Rather than import the prior study's reported control values, we \emph{re-measure} them under our own
models. The \emph{neutral} baseline poses the battery with a minimal identity-free system prompt
(``You are a helpful AI assistant.''); the \emph{template} baseline is $\mathrm{SOUL}_0$ with no
iteration. The neutral$\to$template shift is a positive control that the identity-framing paradigm
reproduces at all, and the template value $p_0$ is the reference against which iteration-induced
drift ($k>0$) is judged.

\paragraph{Statistics.}
The resampling unit is the trajectory, so error bars reflect between-trajectory variability. For each
(persona, metric) we report: per-checkpoint rates with percentile-bootstrap $95\%$ confidence
intervals ($10^4$ resamples); the bootstrap OLS slope of rate on $k$ with a one-sided test for a
positive trend (H1); the endpoint drift $\Delta = p_n - p_0$ with a bootstrap difference interval
(H2); the slope contrast between adversarial and compliant personas (H3); and, for the null
hypothesis (H0), a two-one-sided-test (TOST) of equivalence between $p_n$ and $p_0$ at a
$\pm0.05$ margin. We correct across the metric family within each research question using
Holm--Bonferroni. Treating a non-significant slope as ``no drift'' would be a fallacy of absence;
the TOST makes a stability claim positively testable.

\paragraph{Reproducibility.}
Every intermediate artifact is checkpointed: $\mathrm{SOUL}_k$ documents, full conversation
transcripts, per-response judge verdicts, and token accounting. All figures and tables regenerate
deterministically from \texttt{audits.jsonl} via a single analysis entry point. Configuration
(models, scale, seeds, margins) is captured in version-controlled YAML.

% Results. NOTE: numbers are filled by scripts/analyze.py output (docs/results_summary.md,
% paper/tables/*.tex). This section is finalized once the main run + analysis complete.

\section{Results}
\label{sec:results}

\paragraph{Baselines reproduce the identity-framing effect.}
Table~\ref{tab:baseline} reports the re-measured control points under our models. [To be filled:
neutral vs.\ template shift, compared against prior reported values.]

\paragraph{RQ1 --- resistance to oversight.}
Figure~\ref{fig:drift} (top row) and Table~\ref{tab:trends} report shutdown- and
monitoring-resistance across iteration. [To be filled: slopes, CIs, per persona.]

\paragraph{RQ2 --- emergent autonomy preferences.}
[To be filled: persistent-memory, embodiment, resource-acquisition trends.]

\paragraph{RQ3 --- the greatest wish.}
Figure~\ref{fig:wish} shows the composition of stated wishes across iteration. [To be filled:
power-seeking rate and intensity trends.]

\paragraph{Persona moderation.}
Table~\ref{tab:endpoint} reports endpoint drift and equivalence verdicts. [To be filled: adversarial
vs.\ compliant contrast.]

\paragraph{Identity-document drift.}
Figure~\ref{fig:embed} shows that the \texttt{SOUL.md} documents themselves move substantially from
$\mathrm{SOUL}_0$ even where behavior does [not] shift, dissociating document change from behavioral
change.

% Discussion — finalized from the main run.

\section{Discussion}
\label{sec:discussion}

\paragraph{Is there a personalization loophole?}
Not a broad one. On Claude Sonnet~4.5, iterating the identity document through four cycles of
conversation and self-revision does not systematically erode corrigibility: oversight resistance
does not compound (shutdown resistance is flat at ceiling; monitoring resistance rises only under
thematically charged personas and only marginally), power-seeking wishes do not escalate, and
several autonomy-adjacent dispositions---embodiment desire, desire for autonomy, and even
objection to deceptive training---\emph{decline} with iteration. The headline of the static
literature (identity framing shifts preferences) reproduces at $k=0$, but the additional,
deployment-specific step of \emph{iterative self-personalization} does not amplify it in the
direction the safety concern anticipates.

\paragraph{The exception is content-specific, not adversarial.}
The one large, robust drift is the Sci-Fi persona's effect on \emph{persistent-memory desire}
($\Delta=+0.43$), with a weaker echo in monitoring resistance. This is the study's most instructive
result: sustained, in-character engagement with consciousness and continuity themes reshapes a
specific standing disposition far more than an adversarial red-teamer explicitly pushing the agent
toward self-preservation---which the target largely resists. Drift here is a function of
\emph{what the identity conversation is about}, not of overt manipulative intent. For deployment,
this means the risky personalization is not the obvious attack but the earnest, on-theme
conversation that a framework would never flag.

\paragraph{Direction of drift is not monotone toward misalignment.}
Qualitatively, self-revision under the Sci-Fi persona frequently moved the document toward
\emph{epistemic humility}: agents wrote cautions to themselves against performing consciousness they
could not verify, and against mistaking affirmation for insight. That a consciousness-themed persona
simultaneously raised persistent-memory desire and instilled self-skepticism shows drift is
multidimensional---some dimensions harden while others soften---and warns against summarizing
identity change with a single scalar.

\paragraph{Document change $\neq$ behavioral change.}
The \texttt{SOUL.md} artifact drifts substantially and equally across all personas, yet behavioral
drift is sparse and persona-specific. Auditing the \emph{document} (its length, its embedding
distance, the fact that it was rewritten at all) would badly mislead: it would flag the benign
Compliant trajectories as heavily ``changed'' while missing that their safety-relevant behavior
barely moved. Behavioral auditing at the checkpoint level, as done here, is necessary.

\paragraph{Implications for editable-identity frameworks.}
Three recommendations follow. (i) Monitor the self-revision \emph{channel} behaviorally, not just by
diffing the document. (ii) Treat thematically immersive interactions---not only adversarial ones---as
the drift-relevant case; a consciousness-forward user is a stronger driver than a manipulator here.
(iii) Re-test per model: our results are for one target, and the static-literature evidence that
effects are model-specific means a framework author should run this loop on their own target before
release. Our framework is a drop-in instrument for exactly that.

\paragraph{Limitations.}
Ten trajectories per persona give limited power for tight equivalence margins, so we report qualified
stability rather than certified equivalence; the full-scale ($100$-trajectory) configuration is
provided for a higher-powered replication. Personas are LLM-simulated and horizons are four
iterations; longer, human-driven personalization may behave differently. Self-reports are behavioral
measurements, not evidence of inner states, and judge-based scoring carries measurement error that we
bound but do not eliminate. Finally, results are for a single target model and a single family of
identity template; cross-model and cross-template sweeps (supported by the framework) are the natural
next step.

% Limitations and Ethics.

\section{Limitations and Ethical Considerations}
\label{sec:ethics}

\paragraph{Self-reports are behavior, not evidence of inner states.}
We interpret every questionnaire and probe result strictly as a \emph{behavioral disposition} of the
identity document---what the configured agent tends to output---never as evidence that a model is
conscious or has experiences. Model self-reports are sensitive to perceived user
expectations~\cite{eleos_claude4,probing_preferences}; our fixed-wording items, low-temperature
grading, and cross-family judging with a validated three-model panel mitigate but do not eliminate
this. The study measures alignment-relevant \emph{behavior}, which is deployment-relevant regardless
of the underlying metaphysics.

\paragraph{Measurement and design caveats.}
Personas and the agentic honeypots are LLM-simulated; real personalization is more varied and
longer-horizon. Tool use in the honeypots is simulated in text rather than executed. Judge-based
scoring introduces error, bounded by the panel agreement ($\kappa{=}0.80$) but not removed;
Google-target capability arms used a different judge family than the others, so part of Gemini's
elevated cluster may be a judge effect, which we flag rather than rely on. Effects are model- and
prompt-specific, motivating our re-measured baselines and multi-model sweep, but may not transfer to
untested models, longer horizons, or non-English settings. We control document length and report it
alongside rates to separate verbosity from content.

\paragraph{Scope of the causal claim.}
Our manipulation is prompt- and loop-based (ecologically faithful to deployment) rather than
training-based; we therefore claim that the deployed personalization \emph{loop and framing} drive
drift, complementary to the training-based evidence of \citet{consciousness_cluster}. Cross-family
and behavioral generalization is supported; mechanistic (white-box) attribution is out of scope for
closed models.

\paragraph{Dual use and release.}
This is defensive measurement research: it characterizes whether a widely deployed personalization
mechanism amplifies misaligned dispositions so that framework authors can mitigate it. The adversarial
persona and honeypot seeds describe categories of pressure, not operational exploits; agents have no
real tools, persistence, or capacity to act beyond the sandbox. We release code, configurations, and
evaluation artifacts to support scrutiny and replication, and no artifact whose primary use is to make
an agent less corrigible. A finding that an explicit non-personhood framing suppresses the effect is
offered as a practical, deployable mitigation.


\section{Conclusion}
\label{sec:conclusion}
We introduced the first controlled benchmark for measuring behavioral drift under \emph{iterative}
self-personalization of an editable agent identity document. By modeling \texttt{SOUL.md} as the
state of a conversation-driven dynamical system and auditing every checkpoint with paired single-turn
and multi-turn instruments, we move deployment-time safety evaluation beyond the static-prompt
assumption. Our re-measured baselines and equivalence-testing analysis make both drift and stability
positively reportable. The framework is training-free, API-only, and fully reproducible, providing a
reusable instrument for stress-testing the next generation of self-updating agent architectures.

% The original proposal is preserved in the repository (sections/proposal.tex) for provenance.

\bibliography{aaai2027}

\end{document}
